\clearpage\thispagestyle{empty}
\pdfbookmark[0]{Quellen zu Komet (Dezember 1811, Monatliche Correspondenz)}{comet1811decmc-sources}
\begin{center}\textbf{Quellen zu Komet (Dezember 1811, Monatliche Correspondenz)}\end{center}
\noindent Diese KI-Übertragung ins heutige Deutsch folgt den \href{https://gdz.sub.uni-goettingen.de/id/PPN236007467}{\textit{Werken}, Band~VI} (Göttingen, 1874), S.~343--344. Der Erstdruck erschien als Teil von Nr.~LVII, \textit{IIter Comet vom Jahre 1811}, in der \textit{Monatlichen Correspondenz}, Band~XXIV, Dezember 1811, S.~\href{https://api.digitale-sammlungen.de/iiif/image/v2/bsb10538613_00611/full/1600,/0/default.jpg}{595}, \href{https://api.digitale-sammlungen.de/iiif/image/v2/bsb10538613_00612/full/1600,/0/default.jpg}{596}, \href{https://api.digitale-sammlungen.de/iiif/image/v2/bsb10538613_00613/full/1600,/0/default.jpg}{597} und \href{https://api.digitale-sammlungen.de/iiif/image/v2/bsb10538613_00614/full/1600,/0/default.jpg}{598}. Der Göttinger Auszug beginnt auf S.~596 nach den Berichten über Seeberg und den Beobachtungen von Olbers und Blanpain. Das digitalisierte Exemplar stammt aus der Bayerischen Staatsbibliothek.
\par\smallskip\noindent Die Werkausgabe setzt das Datum von Gauß’ Brief in den laufenden Text. Die Zeitschrift hatte es in einer Fußnote angegeben. Zahlen, Zeichensetzung in den Tabellen, die Form Beoobachter und beide öffnenden Anführungszeichen folgen der Werkausgabe. Für den 11. Dezember stehen hier bei der Rektaszension 18.0 und bei der Deklination 46.4 Sekunden. Der gesonderte Dezemberbericht in den \textit{Göttingischen gelehrten Anzeigen} gibt 20.5 und 54.7 an. KI-Übertragung und KI-Satz.
\clearpage
